% ===========================================================================
%  Chapitre 1 — Limites et continuité
%  PE 2026, composante ANALYSE
% ===========================================================================
\bmtchapitre{Limites et continuité}
  {Déterminer la limite d'une fonction, y compris d'une composée ; établir la
   continuité d'une fonction en un point et sur un intervalle ; prolonger une
   fonction par continuité ; démontrer qu'une équation admet une solution dans
   un intervalle donné à l'aide du théorème des valeurs intermédiaires.}
  {Analyse \textbullet\ environ 10 heures}

En première, vous avez appris à calculer des limites. Le mot était déjà là,
les formes indéterminées aussi. Ce qui change cette année tient en une
phrase : la limite cesse d'être un calcul pour devenir un argument. On ne
demandera plus seulement \emph{combien vaut} la limite, mais ce qu'elle permet
de \emph{démontrer} — qu'une équation a une solution, qu'une courbe reste
coincée entre deux autres, qu'une fonction se prolonge là où elle n'était pas
définie.

Ce chapitre ouvre l'année et il ouvre aussi tout le reste de l'analyse. Les
deux théorèmes qui le terminent, celui des gendarmes et celui des valeurs
intermédiaires, reviendront sur les intégrales, sur le logarithme, sur les
suites. Ils méritent qu'on s'y arrête.

% ---------------------------------------------------------------------------
\begin{revision}
Rien ici ne concerne le chapitre qui vient : on vérifie seulement que la
première année est encore disponible.

\begin{enumerate}[label=\textbf{\arabic*.}, leftmargin=*, itemsep=0.35em]
  \item Donner l'ensemble de définition de $x \mapsto \dfrac{2x-1}{x^{2}-9}$,
        puis celui de $x \mapsto \sqrt{5-x}$.
  \item Factoriser $x^{2}-5x+6$, puis $x^{3}-8$.
  \item Écrire sans valeur absolue $\abs{x-3}$ selon la position de $x$ par
        rapport à $3$.
  \item Rappeler les limites en $+\infty$ de $x \mapsto x^{n}$ ($n \geq 1$),
        de $x \mapsto \dfrac{1}{x}$ et de $x \mapsto \sqrt{x}$.
  \item Quelles sont les quatre formes indéterminées rencontrées en première ?
  \item Traduire par un intervalle : $\abs{x-2} < 0{,}1$.
\end{enumerate}
\end{revision}

\begin{automatismes}
Sans calculatrice, sans brouillon.

\begin{enumerate}[label=\textbf{\alph*)}, leftmargin=*, itemsep=0.25em]
  \item $\limpinf (3x^{2}-x)$
  \item $\limpinf \dfrac{4x+1}{x}$
  \item $\displaystyle\lim_{x \to 0^{+}} \dfrac{1}{x}$
  \item $\limminf \dfrac{1}{x^{2}}$
  \item $\limx{2} (x^{2}-4)$
  \item $\limpinf \dfrac{5x^{2}-3}{2x^{2}+7}$
  \item $\limx{3} \dfrac{x^{2}-9}{x-3}$
  \item $\limpinf \left(\sqrt{x}-x\right)$
\end{enumerate}
\end{automatismes}

\begin{corrige}[automatismes]
a) $+\infty$ \quad b) $4$ \quad c) $+\infty$ \quad d) $0$ \quad e) $0$ \quad
f) $\tfrac{5}{2}$ \quad g) $6$ (on factorise avant de conclure) \quad
h) $-\infty$ (le terme $-x$ l'emporte).
\end{corrige}

% ===========================================================================
\section{La limite d'une fonction composée}

Les opérations sur les limites vues en première couvrent la somme, le produit
et le quotient. Il manque une pièce, et c'est la plus utile : que devient
$g\bigl(f(x)\bigr)$ quand on connaît séparément la limite de $f$ et celle de
$g$ ?

\begin{theoreme}[limite d'une composée]
Soient $a$, $b$ et $\ell$ trois éléments de $\R \cup \{-\infty ; +\infty\}$.
Si
\[
  \limx{a} f(x) = b
  \qquad\text{et}\qquad
  \lim_{X \to b} g(X) = \ell,
\]
alors la fonction composée $g \circ f$ admet une limite en $a$ et
\[
  \limx{a} g\bigl(f(x)\bigr) = \ell .
\]
\end{theoreme}

\begin{demonstration}[dans le cas où $a$, $b$ et $\ell$ sont réels]
Soit $\varepsilon > 0$. Puisque $g$ tend vers $\ell$ en $b$, il existe
$\alpha > 0$ tel que
$\abs{X - b} < \alpha \implies \abs{g(X) - \ell} < \varepsilon$.
Ce nombre $\alpha$ étant fixé, puisque $f$ tend vers $b$ en $a$, il existe
$\eta > 0$ tel que $\abs{x - a} < \eta \implies \abs{f(x) - b} < \alpha$.
En enchaînant les deux implications : dès que $\abs{x-a} < \eta$, on a
$\abs{f(x) - b} < \alpha$, donc $\abs{g(f(x)) - \ell} < \varepsilon$.
C'est exactement la définition de $\limx{a} g(f(x)) = \ell$.
\end{demonstration}

\bmtmarge{On lit toujours une composée de l'intérieur vers l'extérieur : on
détermine d'abord vers quoi tend ce qui est sous le radical, puis ce que
devient le radical.}

\begin{exemple}[une racine dont l'intérieur tend vers l'infini]
Calculons $\limpinf \sqrt{\dfrac{3x^{2}+1}{x^{2}+2}}$.

On pose $f(x) = \dfrac{3x^{2}+1}{x^{2}+2}$ et $g(X) = \sqrt{X}$, de sorte que
la fonction étudiée est $g \circ f$.

\smallskip
\emph{Limite intérieure.} Les deux termes dominants sont $3x^{2}$ au numérateur
et $x^{2}$ au dénominateur, donc
$\limpinf f(x) = \limpinf \dfrac{3x^{2}}{x^{2}} = 3$.

\smallskip
\emph{Limite extérieure.} La fonction racine est continue en $3$, donc
$\lim_{X \to 3} \sqrt{X} = \sqrt{3}$.

\smallskip
\emph{Conclusion.} Par composition,
$\limpinf \sqrt{\dfrac{3x^{2}+1}{x^{2}+2}} = \sqrt{3}$.
\end{exemple}

\begin{entrainement}[limites de composées]
\exo[\niveau{1}] Déterminer $\limpinf \sqrt{4x^{2}+x}$ et
$\limpinf \left(\dfrac{1}{x+1}\right)^{3}$.

\exo[\niveau{2}] Déterminer $\displaystyle\lim_{x \to 2} \sqrt{\dfrac{x^{2}-4}{x-2}}$.
On commencera par simplifier l'expression sous le radical.

\exo[\niveau{2}] Soit $f(x) = \dfrac{2x-3}{x+1}$. Déterminer les limites de
$\sqrt{f(x)}$ en $+\infty$ et en $-\infty$, après avoir vérifié que
l'expression a bien un sens au voisinage de ces deux infinis.
\end{entrainement}

\begin{corrige}
\textbf{1.} $\limpinf \sqrt{4x^{2}+x} = +\infty$ car $4x^{2}+x \to +\infty$.
Et $\dfrac{1}{x+1} \to 0$, donc son cube tend vers $0$.

\textbf{2.} Pour $x \neq 2$, $\dfrac{x^{2}-4}{x-2} = x+2$, qui tend vers $4$
quand $x \to 2$. Par composition, la limite vaut $\sqrt{4} = 2$.

\textbf{3.} $f(x) \to 2$ aux deux infinis. Il reste à vérifier la positivité :
$f(x) > 0$ dès que $x > \tfrac{3}{2}$ ou $x < -1$, ce qui est le cas au
voisinage de $\pm\infty$. La limite vaut donc $\sqrt{2}$ dans les deux cas.
\end{corrige}

% ===========================================================================
\section{Encadrer pour conclure : comparaison et gendarmes}

Il arrive qu'une fonction résiste au calcul direct — typiquement lorsqu'elle
contient un sinus ou un cosinus, qui n'ont pas de limite en l'infini. On
renonce alors à la calculer et on la coince entre deux fonctions plus simples.

\begin{theoreme}[théorème de comparaison]
Soit $a \in \R \cup \{-\infty ; +\infty\}$ et soient $f$ et $g$ deux fonctions
telles que $f(x) \leq g(x)$ au voisinage de $a$.
\begin{itemize}[leftmargin=1.4em, itemsep=0.2em]
  \item Si $\limx{a} f(x) = +\infty$, alors $\limx{a} g(x) = +\infty$.
  \item Si $\limx{a} g(x) = -\infty$, alors $\limx{a} f(x) = -\infty$.
\end{itemize}
\end{theoreme}

\begin{theoreme}[théorème des gendarmes]
Soient $u$, $f$ et $v$ trois fonctions telles que, au voisinage de $a$,
\[
  u(x) \leq f(x) \leq v(x).
\]
Si $u$ et $v$ admettent la \emph{même} limite finie $\ell$ en $a$, alors $f$
admet aussi $\ell$ pour limite en $a$.
\end{theoreme}

\begin{demonstration}
Soit $\varepsilon > 0$. Comme $u \to \ell$, il existe un voisinage de $a$ sur
lequel $\ell - \varepsilon < u(x)$. Comme $v \to \ell$, il existe un voisinage
de $a$ sur lequel $v(x) < \ell + \varepsilon$. Sur l'intersection de ces deux
voisinages, l'encadrement $u(x) \leq f(x) \leq v(x)$ donne
\[
  \ell - \varepsilon < u(x) \leq f(x) \leq v(x) < \ell + \varepsilon,
\]
c'est-à-dire $\abs{f(x) - \ell} < \varepsilon$. La limite de $f$ en $a$ est
donc $\ell$.
\end{demonstration}

\begin{visualisation}[deux gendarmes qui se rejoignent]
\begin{center}
\begin{tikzpicture}[scale=1]
\begin{axis}[
  width = 11.4cm, height = 6.2cm,
  axis lines = middle,
  xlabel = {$x$}, ylabel = {},
  xmin = 0.55, xmax = 9.6, ymin = -0.85, ymax = 2.5,
  xtick = \empty, ytick = {1}, yticklabels = {$\ell$},
  every axis y label/.style = {at={(ticklabel* cs:1.02)}, anchor=south},
  clip = false, axis line style = {bmtGris},
  tick label style = {font=\footnotesize, color=bmtGris},
]
  \addplot[bmtGris, dashed, domain=0.7:9.4, samples=2] {1};
  \addplot[bmtVetiver, line width=1pt, domain=0.9:9.4, samples=200]
    {1 + 1.6/x};
  \addplot[bmtVetiver, line width=1pt, domain=0.9:9.4, samples=200]
    {1 - 1.6/x};
  \addplot[bmtLaterite, line width=1.3pt, domain=0.9:9.4, samples=300]
    {1 + 1.6*sin(deg(4*x))/x};
  \node[bmtVetiver, font=\footnotesize, anchor=west] at (axis cs:9.45,1.28) {$v$};
  \node[bmtVetiver, font=\footnotesize, anchor=west] at (axis cs:9.45,0.72) {$u$};
  \node[bmtLaterite, font=\footnotesize, anchor=west] at (axis cs:9.45,1.02) {$f$};
\end{axis}
\end{tikzpicture}
\end{center}

Regardez les deux courbes vertes : elles se referment l'une sur l'autre en
allant vers la droite. La courbe rouge oscille entre les deux, sans jamais
pouvoir en sortir. Elle n'a pas le choix : puisque les deux gendarmes
convergent vers $\ell$, elle converge vers $\ell$ aussi. Notez bien qu'on n'a
jamais eu besoin de savoir ce que fait $f$ dans le détail.
\end{visualisation}

\begin{exemple}[la limite qu'on ne peut pas calculer directement]
Déterminons $\limpinf \dfrac{\sin x}{x}$.

La fonction sinus n'a pas de limite en $+\infty$ : elle oscille indéfiniment
entre $-1$ et $1$. Le calcul direct est donc sans issue. En revanche, cet
encadrement est vrai pour tout réel $x$ :
\[
  -1 \leq \sin x \leq 1 .
\]
Pour $x > 0$, on divise les trois membres par $x$ sans changer le sens des
inégalités :
\[
  -\frac{1}{x} \leq \frac{\sin x}{x} \leq \frac{1}{x}.
\]
Or $\limpinf \left(-\dfrac{1}{x}\right) = 0$ et
$\limpinf \dfrac{1}{x} = 0$. Les deux gendarmes ont la même limite, donc
\[
  \limpinf \frac{\sin x}{x} = 0 .
\]
\end{exemple}

\begin{erreurclassique}
Le théorème des gendarmes exige que les deux fonctions encadrantes aient la
\textbf{même} limite. Écrire « $u(x) \leq f(x) \leq v(x)$ avec $u \to 2$ et
$v \to 5$, donc $f \to$ quelque chose entre $2$ et $5$ » ne démontre rien du
tout : $f$ peut très bien n'avoir aucune limite. Quand les deux bornes
diffèrent, l'encadrement ne donne qu'un encadrement, pas une limite.
\end{erreurclassique}

\begin{entrainement}[comparaison et encadrement]
\exo[\niveau{1}] Montrer que $\limpinf (x + \cos x) = +\infty$.

\exo[\niveau{2}] Déterminer $\limpinf \dfrac{2x + \sin x}{x}$.

\exo[\niveau{2}] Déterminer $\displaystyle\lim_{x \to 0} x^{2}\sin\!\left(\dfrac{1}{x}\right)$.

\exo[\niveau{3}] Soit $f$ une fonction vérifiant, pour tout $x \geq 1$,
$\abs{f(x) - 3} \leq \dfrac{1}{x^{2}}$. Déterminer $\limpinf f(x)$ en
justifiant par un encadrement.
\end{entrainement}

\begin{corrige}
\textbf{1.} Pour tout $x$, $\cos x \geq -1$, donc $x + \cos x \geq x - 1$.
Comme $x - 1 \to +\infty$, le théorème de comparaison conclut.

\textbf{2.} $\dfrac{2x+\sin x}{x} = 2 + \dfrac{\sin x}{x}$, et le second terme
tend vers $0$ (exemple ci-dessus) : la limite vaut $2$.

\textbf{3.} Pour $x \neq 0$, $-x^{2} \leq x^{2}\sin\!\left(\tfrac{1}{x}\right) \leq x^{2}$,
et les deux bornes tendent vers $0$ : la limite vaut $0$.

\textbf{4.} L'hypothèse s'écrit $3 - \dfrac{1}{x^{2}} \leq f(x) \leq 3 + \dfrac{1}{x^{2}}$.
Les deux bornes tendent vers $3$, donc $\limpinf f(x) = 3$.
\end{corrige}

% ===========================================================================
\section{La continuité}

\begin{definition}[continuité en un point]
Soit $f$ une fonction définie sur un intervalle $I$ et soit $a \in I$. On dit
que $f$ est \textbf{continue en $a$} lorsque
\[
  \limx{a} f(x) = f(a).
\]
Lorsque $f$ est continue en tout point de $I$, on dit que $f$ est
\textbf{continue sur $I$}.
\end{definition}

\bmtmarge{Trois conditions se cachent dans cette égalité : $f$ doit être
définie en $a$, la limite doit exister, et les deux nombres doivent
coïncider.}

\begin{propriete}[fonctions continues usuelles]
Sont continues sur tout intervalle contenu dans leur ensemble de définition :
les fonctions polynômes, les fonctions rationnelles, la fonction racine
carrée, la fonction valeur absolue, les fonctions sinus et cosinus. De plus,
toute somme, tout produit, tout quotient (à dénominateur non nul) et toute
composée de fonctions continues est continue.
\end{propriete}

\begin{visualisation}[là où la courbe se casse]
\begin{center}
\begin{tikzpicture}
\begin{axis}[
  width = 11.4cm, height = 5.6cm,
  axis lines = middle,
  xlabel = {$x$},
  xmin = -0.6, xmax = 4.4, ymin = -0.6, ymax = 3.3,
  xtick = {2}, ytick = {1,2}, xticklabels = {$a$},
  axis line style = {bmtGris},
  tick label style = {font=\footnotesize, color=bmtGris},
  clip = false,
]
  \addplot[bmtIndigo, line width=1.2pt, domain=-0.4:2, samples=60] {0.5*x + 0.4};
  \addplot[bmtIndigo, line width=1.2pt, domain=2:4.2, samples=60] {0.35*x + 1.6};
  \filldraw[bmtIndigo] (axis cs:2,1.4) circle (2.1pt);
  \draw[bmtIndigo, line width=1pt, fill=bmtPapier] (axis cs:2,2.3) circle (2.1pt);
  \draw[bmtLaterite, dashed, line width=0.8pt] (axis cs:2,1.4) -- (axis cs:2,2.3);
  \node[bmtLaterite, font=\footnotesize, anchor=west] at (axis cs:2.12,1.85) {le saut};
\end{axis}
\end{tikzpicture}
\end{center}

Suivez la courbe de la gauche vers la droite : arrivé en $a$, il faut lever le
crayon pour continuer. La fonction est bien définie en $a$ — le point plein
indique sa valeur — mais la limite à droite ne redescend pas jusqu'à lui.
C'est précisément ce que la définition interdit. Une fonction continue sur un
intervalle, c'est une courbe qu'on trace sans lever le crayon ; retenez
l'image, mais souvenez-vous que la démonstration, elle, passe par la limite.
\end{visualisation}

\subsection{Prolonger une fonction par continuité}

Une fonction peut être parfaitement régulière partout sauf en un point où elle
n'est simplement pas définie — un dénominateur qui s'annule, alors même que le
numérateur s'annule aussi. Le trou est alors réparable.

\begin{definition}[prolongement par continuité]
Soit $f$ définie sur $I \setminus \{a\}$. Si $f$ admet une limite finie $\ell$
en $a$, alors la fonction $\tilde{f}$ définie par
\[
  \tilde{f}(x) =
  \begin{cases}
    f(x) & \text{si } x \in I \setminus \{a\}, \\[0.3em]
    \ell & \text{si } x = a,
  \end{cases}
\]
est continue en $a$. On l'appelle le \textbf{prolongement par continuité} de
$f$ en $a$.
\end{definition}

\begin{exemple}[réparer un trou]
Soit $f(x) = \dfrac{x^{2}-9}{x-3}$, définie sur $\R \setminus \{3\}$.

Pour $x \neq 3$, on factorise le numérateur :
$\dfrac{x^{2}-9}{x-3} = \dfrac{(x-3)(x+3)}{x-3} = x+3$.
Donc $\limx{3} f(x) = 6$, limite finie.

La fonction $\tilde{f}$ définie sur $\R$ par $\tilde{f}(x) = f(x)$ pour
$x \neq 3$ et $\tilde{f}(3) = 6$ est le prolongement par continuité de $f$
en $3$. On reconnaît d'ailleurs $\tilde{f}(x) = x+3$ pour tout $x$.
\end{exemple}

\begin{avertissement}
Le prolongement n'est possible que si la limite est \emph{finie}. La fonction
$x \mapsto \dfrac{1}{x}$ ne se prolonge pas en $0$ : sa limite y est infinie,
et aucune valeur réelle attribuée à $\tilde{f}(0)$ ne rendra la fonction
continue.
\end{avertissement}

\begin{entrainement}[continuité et prolongement]
\exo[\niveau{1}] La fonction $f$ définie par $f(x) = 2x - 1$ si $x \leq 1$ et
$f(x) = x^{2}$ si $x > 1$ est-elle continue en $1$ ?

\exo[\niveau{2}] Étudier la possibilité de prolonger par continuité en $2$ la
fonction $x \mapsto \dfrac{x^{2}-4}{x^{2}-3x+2}$.

\exo[\niveau{3}] Déterminer le réel $m$ pour que la fonction définie par
$f(x) = \dfrac{\sqrt{x+1}-1}{x}$ si $x \neq 0$ et $f(0) = m$ soit continue
en $0$.
\end{entrainement}

\begin{corrige}
\textbf{1.} À gauche : $\lim_{x \to 1^{-}} (2x-1) = 1$. À droite :
$\lim_{x \to 1^{+}} x^{2} = 1$. Et $f(1) = 1$. Les trois nombres coïncident,
donc $f$ est continue en $1$.

\textbf{2.} $\dfrac{x^{2}-4}{x^{2}-3x+2} = \dfrac{(x-2)(x+2)}{(x-2)(x-1)} = \dfrac{x+2}{x-1}$
pour $x \neq 2$, expression qui tend vers $4$ quand $x \to 2$. Le prolongement
existe, avec la valeur $4$.

\textbf{3.} On multiplie par la quantité conjuguée :
$\dfrac{\sqrt{x+1}-1}{x} = \dfrac{x}{x\left(\sqrt{x+1}+1\right)} = \dfrac{1}{\sqrt{x+1}+1}$,
qui tend vers $\tfrac{1}{2}$. Donc $m = \tfrac{1}{2}$.
\end{corrige}

% ===========================================================================
\section{Le théorème des valeurs intermédiaires}

Voici le résultat vers lequel tout le chapitre penchait. Il transforme la
continuité, qui est une propriété locale, en une garantie d'existence.

\begin{theoreme}[valeurs intermédiaires]
Soit $f$ une fonction continue sur un intervalle $\intervalle{a}{b}$. Alors,
pour tout réel $k$ compris entre $f(a)$ et $f(b)$, il existe au moins un réel
$c \in \intervalle{a}{b}$ tel que
\[
  f(c) = k .
\]
\end{theoreme}

\begin{visualisation}[toute valeur intermédiaire est atteinte]
\begin{center}
\begin{tikzpicture}
\begin{axis}[
  width = 11.4cm, height = 6cm,
  axis lines = middle,
  xlabel = {$x$},
  xmin = -0.5, xmax = 5.3, ymin = -1.6, ymax = 3.4,
  xtick = {0.5, 4.6}, xticklabels = {$a$, $b$},
  ytick = {-1, 0.9, 2.8}, yticklabels = {$f(a)$, $k$, $f(b)$},
  axis line style = {bmtGris},
  tick label style = {font=\footnotesize, color=bmtGris},
  clip = false,
]
  % La courbe vaut exactement -1 en a et 2,8 en b ; elle coupe la droite
  % y = 0,9 en trois points, dont les abscisses sont calculees ci-dessous.
  \addplot[bmtVetiver, line width=1.3pt, domain=0.5:4.6, samples=250]
    {-1 + 3.8*(x-0.5)/4.1 + 1.5*sin(360*(x-0.5)/4.1)};
  \addplot[bmtLaterite, dashed, line width=0.8pt, domain=0.5:4.6, samples=2] {0.9};
  \draw[bmtGris, dotted] (axis cs:0.5,-1) -- (axis cs:0.5,0);
  \draw[bmtGris, dotted] (axis cs:4.6,2.8) -- (axis cs:4.6,0);
  \filldraw[bmtLaterite] (axis cs:1.168,0.9) circle (1.9pt);
  \filldraw[bmtLaterite] (axis cs:2.550,0.9) circle (1.9pt);
  \filldraw[bmtLaterite] (axis cs:3.932,0.9) circle (1.9pt);
  \node[bmtLaterite, font=\footnotesize, anchor=north] at (axis cs:1.168,0.72) {$c_{1}$};
  \node[bmtLaterite, font=\footnotesize, anchor=north] at (axis cs:2.550,0.72) {$c_{2}$};
  \node[bmtLaterite, font=\footnotesize, anchor=north] at (axis cs:3.932,0.72) {$c_{3}$};
\end{axis}
\end{tikzpicture}
\end{center}

Tracez la droite horizontale d'ordonnée $k$ : puisque la courbe part en
dessous et arrive au-dessus sans jamais se casser, elle doit franchir cette
droite. Le théorème affirme l'existence d'au moins une traversée — ici, il y
en a trois. Retenez donc ce que le théorème ne dit pas : il ne donne ni la
valeur de $c$, ni le nombre de solutions.
\end{visualisation}

\begin{theoreme}[cas d'une fonction strictement monotone]
Si de plus $f$ est \emph{strictement monotone} sur $\intervalle{a}{b}$, alors
le réel $c$ est \textbf{unique}. Autrement dit, l'équation $f(x) = k$ admet
une solution et une seule dans $\intervalle{a}{b}$.
\end{theoreme}

\begin{demonstration}
L'existence est donnée par le théorème précédent. Pour l'unicité, supposons
qu'il existe $c_{1} < c_{2}$ dans $\intervalle{a}{b}$ avec
$f(c_{1}) = f(c_{2}) = k$. Si $f$ est strictement croissante, alors
$c_{1} < c_{2}$ entraîne $f(c_{1}) < f(c_{2})$, ce qui contredit l'égalité.
Le raisonnement est identique pour une fonction strictement décroissante. Il
ne peut donc exister deux tels réels.
\end{demonstration}

\begin{methode}{montrer qu'une équation admet une unique solution}
Face à une équation du type $f(x) = k$ sur un intervalle, la rédaction attendue
tient en quatre points, et le correcteur les cherche un par un.

\begin{enumerate}[label=\textbf{\arabic*.}, leftmargin=*, itemsep=0.3em]
  \item \textbf{Continuité} : dire pourquoi $f$ est continue sur l'intervalle
        (polynôme, quotient à dénominateur non nul, composée de fonctions
        continues\ldots).
  \item \textbf{Monotonie stricte} : l'établir par le signe de la dérivée,
        et le dire explicitement.
  \item \textbf{Encadrement de $k$} : calculer $f(a)$ et $f(b)$, puis vérifier
        que $k$ est bien compris entre les deux.
  \item \textbf{Conclusion} : invoquer le théorème par son nom et annoncer
        l'existence et l'unicité.
\end{enumerate}
\end{methode}

\begin{exemple}[une équation du troisième degré]
Montrons que l'équation $x^{3} + x - 1 = 0$ admet une unique solution dans
$\intervalle{0}{1}$.

Posons $f(x) = x^{3} + x - 1$.

\smallskip
\emph{Continuité.} $f$ est une fonction polynôme, donc continue sur
$\intervalle{0}{1}$.

\smallskip
\emph{Monotonie.} $f'(x) = 3x^{2} + 1 > 0$ pour tout réel $x$, donc $f$ est
strictement croissante sur $\intervalle{0}{1}$.

\smallskip
\emph{Encadrement.} $f(0) = -1$ et $f(1) = 1$. Le nombre $0$ est bien compris
entre $-1$ et $1$.

\smallskip
\emph{Conclusion.} D'après le théorème des valeurs intermédiaires appliqué à
une fonction continue et strictement monotone, l'équation $f(x) = 0$ admet une
solution unique $c$ dans $\intervalle{0}{1}$.

\smallskip
On peut ensuite l'encadrer : $f(0{,}6) \approx -0{,}18 < 0$ et
$f(0{,}7) \approx 0{,}04 > 0$, donc $0{,}6 < c < 0{,}7$.
\end{exemple}

\begin{entrainement}[application du théorème]
\exo[\niveau{1}] Montrer que l'équation $x^{3} = 5$ admet une unique solution
dans $\intervalle{1}{2}$.

\exo[\niveau{2}] Soit $f(x) = x^{3} - 3x + 1$. Montrer que l'équation
$f(x) = 0$ admet exactement trois solutions réelles, et donner un encadrement
d'amplitude $0{,}1$ de chacune.

\exo[\niveau{3}] Soit $f$ continue sur $\intervalle{0}{1}$ à valeurs dans
$\intervalle{0}{1}$. Montrer qu'il existe $c \in \intervalle{0}{1}$ tel que
$f(c) = c$.
\end{entrainement}

\begin{corrige}
\textbf{1.} $f(x) = x^{3}$ est continue et strictement croissante sur
$\intervalle{1}{2}$, avec $f(1) = 1$ et $f(2) = 8$. Comme $1 < 5 < 8$, le
théorème donne l'existence et l'unicité.

\textbf{2.} $f'(x) = 3x^{2} - 3 = 3(x-1)(x+1)$, donc $f$ croît sur
$\left]-\infty ; -1\right]$, décroît sur $\intervalle{-1}{1}$, croît sur
$\left[1 ; +\infty\right[$. On calcule $f(-1) = 3 > 0$ et $f(1) = -1 < 0$.
Sur chacun des trois intervalles de monotonie, $f$ est continue et strictement
monotone et change de signe : le théorème s'applique trois fois. Les solutions
sont approximativement $-1{,}9$, $0{,}3$ et $1{,}5$.

\textbf{3.} $g = f - \mathrm{id}$ est continue sur $\intervalle{0}{1}$ comme
différence de fonctions continues. Or $g(0) = f(0) \geq 0$ et
$g(1) = f(1) - 1 \leq 0$ puisque $f$ est à valeurs dans $\intervalle{0}{1}$.
Le réel $0$ est donc compris entre $g(1)$ et $g(0)$ : il existe
$c \in \intervalle{0}{1}$ tel que $g(c) = 0$, c'est-à-dire $f(c) = c$.
\end{corrige}

\begin{notepourlaclasse}
Le troisième exercice (point fixe) dépasse le programme mais tombe
régulièrement aux concours d'entrée. Il peut être donné en recherche
collective : la difficulté n'est pas technique, elle est de \emph{penser} à
introduire $g$. Laissez cinq minutes de flottement avant de souffler l'indice.

Sur la méthode en quatre points, imposez la rédaction complète dès les
premières copies : c'est là que se perdent le plus de points au baccalauréat,
les élèves concluant l'existence sans jamais justifier la continuité.
\end{notepourlaclasse}


% ===========================================================================
\begin{annales}
Les exercices qui suivent sont repris de sujets réellement posés, à Madagascar
et dans les pays voisins ; leur provenance est indiquée à chaque fois. Tous
n'utilisent que des fonctions polynômes, rationnelles ou irrationnelles :
aucun ne suppose une notion qui n'aurait pas encore été étudiée.

\exoann{Baccalauréat, Côte d'Ivoire, série C, session 2021 — exercice 1}
Répondre par vrai ou faux, en justifiant.
« Si $f$ est croissante et majorée sur $\intervalleo{2}{5}$, alors $f$ admet
pour limite $+\infty$ à gauche en $5$. »

\exoann{Baccalauréat blanc, lycée Philibert Tsiranana, Terminale S, 2022-2023}
Soit $f$ la fonction définie sur $\R \setminus \{2\}$ par
$f(x) = \dfrac{x^{2}-3x+4}{x-2}$.
Déterminer les limites de $f$ en $-\infty$, en $+\infty$, puis à gauche et à
droite de $2$. En déduire l'existence d'une asymptote verticale.

\exoann{Concours d'entrée, ENI Fianarantsoa, session 2002-2003 — exercice 1}
Soit $f_{m}$ la fonction définie par
$f_{m}(x) = \dfrac{x^{2}+3x+4m}{x^{2}+(5m+1)x+3}$, où $m$ est un paramètre
réel, et $(C_{m})$ sa courbe. Pour $m = 0$, déterminer l'ensemble de
définition de $f_{0}$ et ses limites aux bornes.

\exoann{Concours d'entrée, ENI Fianarantsoa, session 2002-2003 — exercice 1, suite}
Avec les mêmes notations, montrer que toutes les courbes $(C_{m})$ passent par
trois points fixes, dont on déterminera les coordonnées.

\exoann{Baccalauréat probatoire, Cameroun, série D, session 2022 — exercice 4}
Soit $f$ la fonction définie sur $\intervallefo{0}{+\infty}$ par
$f(x) = \dfrac{3x}{3+4x}$. Calculer $\lim_{x\to+\infty}f(x)$ et interpréter
graphiquement le résultat.

\exoann{Baccalauréat probatoire, Cameroun, séries D et TI, session 2023 — exercice 4}
Soit $f$ définie par $f(x) = x+\dfrac{1}{x+1}$.
\begin{enumerate}[label=\textbf{\alph*)}, leftmargin=*, itemsep=0.15em]
  \item Déterminer l'ensemble de définition de $f$.
  \item Calculer les limites de $f$ aux bornes de cet ensemble.
  \item En déduire les asymptotes de la courbe de $f$.
\end{enumerate}

\exoann{D'après le baccalauréat probatoire, Cameroun, série C, session 2021 — exercice 3}
Soit $f$ définie sur $\R\setminus\{1\}$ par $f(x) = \dfrac{x^{2}-4}{1-x}$, de
courbe $(C)$. Calculer les limites de $f$ aux bornes de son ensemble de
définition, puis donner l'équation de l'asymptote verticale $(L)$ à $(C)$.

\exoann{D'après le baccalauréat probatoire, Cameroun, série C, session 2023 — exercice 4}
Soit $f$ définie par $f(x) = \dfrac{2x^{2}-6x+3}{2x-3}$. Déterminer son
ensemble de définition, puis ses limites aux bornes de cet ensemble. Qu'en
déduit-on pour la courbe de $f$ ?

\exoann{D'après le baccalauréat probatoire, Cameroun, série C, session 2022 — exercice 2}
Soit $f$ définie sur $\R$ par
$f(x) = \dfrac{1}{3}x^{3}+\dfrac{3}{2}x^{2}-4x-\dfrac{31}{6}$.
\begin{enumerate}[label=\textbf{\alph*)}, leftmargin=*, itemsep=0.15em]
  \item Vérifier que la courbe de $f$ passe par le point $(-1\,;0)$.
  \item Démontrer que l'équation $f(x) = 0$ admet exactement trois solutions
        réelles.
\end{enumerate}

\medskip
\bmtsource{Les autres énoncés d'annales portant sur les limites et la
continuité font tous intervenir le logarithme ou l'exponentielle. Ils
figurent donc aux chapitres 5 et 6, où ils sont signalés par la mention
« limites et continuité ».}
\end{annales}

\begin{corrige}[annales]
\textbf{1.} Faux. Une fonction croissante et majorée sur un intervalle admet
en la borne droite une limite \emph{finie} — la borne supérieure de ses
valeurs. Une limite infinie contredirait le fait qu'elle soit majorée.

\textbf{2.} Aux deux infinis $f(x) \sim x$, donc $f \to \pm\infty$. En $2^{-}$ :
le numérateur tend vers $2>0$ et le dénominateur vers $0^{-}$, donc
$f \to -\infty$ ; en $2^{+}$, $f \to +\infty$. La droite $x = 2$ est asymptote
verticale.

\textbf{3.} Pour $m = 0$ le dénominateur $x^{2}+x+3$ a pour discriminant
$-11<0$ et ne s'annule jamais : $\Df = \R$. Les degrés du numérateur et du
dénominateur étant égaux, la limite en $\pm\infty$ est le quotient des
coefficients dominants, soit $1$. La droite $y = 1$ est asymptote horizontale
dans les deux directions.

\textbf{4.} Un point $(x_{0}\,;y_{0})$ appartient à toutes les courbes si et
seulement si, pour tout $m$,
\[
  x_{0}^{2}+3x_{0}+4m
  = y_{0}\left(x_{0}^{2}+(5m+1)x_{0}+3\right).
\]
En regroupant les termes en $m$ et en annulant séparément les deux paquets, il
vient $5x_{0}y_{0} = 4$ et $x_{0}^{2}+3x_{0} = y_{0}\left(x_{0}^{2}+x_{0}
+3\right)$. En reportant $y_{0} = \frac{4}{5x_{0}}$ dans la seconde égalité,
on obtient $5x_{0}^{3}+11x_{0}^{2}-4x_{0}-12 = 0$, dont $1$ est racine
évidente ; la factorisation $(x_{0}-1)\left(5x_{0}^{2}+16x_{0}+12\right)$
donne les deux autres, $-\frac65$ et $-2$. Les trois points fixes sont donc
\[
  \left(1\,;\tfrac45\right), \qquad
  \left(-\tfrac65\,;-\tfrac23\right), \qquad
  \left(-2\,;-\tfrac25\right).
\]

\textbf{5.} Numérateur et dénominateur sont de degré $1$ : la limite est le
quotient des coefficients de $x$, soit $\frac34$. La droite d'équation
$y = \frac34$ est asymptote horizontale à la courbe en $+\infty$.

\textbf{6. a)} Le dénominateur $x+1$ s'annule en $-1$ :
$\Df = \R\setminus\{-1\}$.
\textbf{b)} En $\pm\infty$, $\frac{1}{x+1}\to0$ et $f(x)\to\pm\infty$. En
$-1^{-}$, $\frac{1}{x+1}\to-\infty$ donc $f\to-\infty$ ; en $-1^{+}$,
$f\to+\infty$.
\textbf{c)} La droite $x = -1$ est asymptote verticale. De plus
$f(x)-x = \frac{1}{x+1}\to0$ en $\pm\infty$ : la droite $y = x$ est asymptote
oblique dans les deux directions.

\textbf{7.} L'ensemble de définition est $\R\setminus\{1\}$. Aux infinis, le
numérateur est de degré $2$ et le dénominateur de degré $1$ : en gardant les
termes dominants, $f(x)$ se comporte comme $\frac{x^{2}}{-x} = -x$, donc
$f(x)\to-\infty$ en $+\infty$ et $f(x)\to+\infty$ en $-\infty$. En $1$, le
numérateur tend vers $-3$, strictement négatif, tandis que $1-x$ tend vers
$0^{+}$ à gauche de $1$ et vers $0^{-}$ à droite : ainsi
$f(x)\to-\infty$ quand $x\to1^{-}$ et $f(x)\to+\infty$ quand $x\to1^{+}$. La
droite $(L)$ d'équation $x = 1$ est asymptote verticale à $(C)$.

\textbf{8.} Le dénominateur s'annule pour $x = \frac32$ :
$\Df = \R\setminus\left\{\frac32\right\}$. Aux infinis, $f(x)$ se comporte
comme $\frac{2x^{2}}{2x} = x$ et tend donc vers $\pm\infty$. En $\frac32$, le
numérateur vaut $2\times\frac94-9+3 = -\frac32$, strictement négatif, et
$2x-3$ tend vers $0^{-}$ à gauche, vers $0^{+}$ à droite :
$f(x)\to+\infty$ quand $x\to\frac32^{-}$, et $f(x)\to-\infty$ quand
$x\to\frac32^{+}$. La droite $x = \frac32$ est asymptote verticale.

\textbf{9. a)} $f(-1) = -\dfrac13+\dfrac32+4-\dfrac{31}{6}
= \dfrac{-2+9+24-31}{6} = 0$ : le point $(-1\,;0)$ appartient bien à la
courbe.
\textbf{b)} $f'(x) = x^{2}+3x-4 = (x+4)(x-1)$, positive à l'extérieur des
racines et négative entre elles. La fonction croît donc sur
$\left]-\infty\,;-4\right]$, décroît sur $\intervalle{-4}{1}$, croît sur
$\intervallefo{1}{+\infty}$, avec
\[
  f(-4) = -\frac{64}{3}+24+16-\frac{31}{6} = \frac{-128+144+96-31}{6}
  = \frac{27}{2} > 0,
  \qquad
  f(1) = \frac{2+9-24-31}{6} = -\frac{22}{3} < 0 .
\]
Sur chacun des trois intervalles, $f$ est continue et strictement monotone, et
elle change de signe : $f$ va de $-\infty$ à $\frac{27}{2}$, puis de
$\frac{27}{2}$ à $-\frac{22}{3}$, puis de $-\frac{22}{3}$ à $+\infty$. Le
théorème des valeurs intermédiaires donne une racine, et une seule, dans
chacun d'eux : l'équation $f(x) = 0$ a exactement trois solutions réelles,
dont $-1$, qui se trouve dans le deuxième intervalle.
\end{corrige}

% ===========================================================================
\begin{jecherche}[la population d'une commune]
Une commune rurale compte $8\,000$ habitants en 2020. Une étude modélise sa
population, en milliers d'habitants, $t$ années après 2020, par
\[
  P(t) = \frac{8 + 3t}{1 + 0{,}2t}, \qquad t \geq 0 .
\]

\begin{enumerate}[label=\textbf{\arabic*.}, leftmargin=*, itemsep=0.3em]
  \item Vérifier que le modèle donne bien $8\,000$ habitants en 2020.
  \item Déterminer $\lim_{t \to +\infty} P(t)$ et interpréter ce nombre pour
        la commune.
  \item Montrer que $P$ est strictement croissante sur $\left[0 ; +\infty\right[$.
  \item En quelle année la population dépassera-t-elle $12\,000$ habitants ?
        Justifier l'existence et l'unicité de cette année avant de la
        déterminer.
  \item Le modèle prévoit-il que la commune atteindra un jour $16\,000$
        habitants ? Que peut-on en dire ?
\end{enumerate}
\end{jecherche}

\begin{corrige}[la population d'une commune]
\textbf{1.} $P(0) = 8$, soit $8\,000$ habitants.

\textbf{2.} $P(t) = \dfrac{8+3t}{1+0{,}2t} \sim \dfrac{3t}{0{,}2t} = 15$ quand
$t \to +\infty$. La population tend vers $15\,000$ habitants sans jamais les
atteindre : le modèle prévoit une saturation.

\textbf{3.} $P'(t) = \dfrac{3(1+0{,}2t) - 0{,}2(8+3t)}{(1+0{,}2t)^{2}}
= \dfrac{3 - 1{,}6}{(1+0{,}2t)^{2}} = \dfrac{1{,}4}{(1+0{,}2t)^{2}} > 0$.

\textbf{4.} $P$ est continue (quotient de fonctions polynômes à dénominateur
non nul sur $\left[0 ; +\infty\right[$) et strictement croissante, avec
$P(0) = 8$ et $\lim_{+\infty} P = 15$. Comme $8 < 12 < 15$, il existe un
unique $t_{0}$ tel que $P(t_{0}) = 12$. On résout :
$8 + 3t = 12 + 2{,}4t$, d'où $t_{0} = \dfrac{4}{0{,}6} \approx 6{,}67$. Le seuil
est franchi au cours de la septième année, donc en 2027.

\textbf{5.} Non : $P$ est majorée par $15$, donc l'équation $P(t) = 16$ n'a
pas de solution. C'est une limite du modèle, pas une prédiction sur la
commune — un modèle rationnel de ce type impose une saturation qui n'est
qu'une hypothèse.
\end{corrige}

% ===========================================================================
\begin{jemeteste}
Répondre par vrai ou faux, en justifiant en une ligne.

\begin{enumerate}[label=\textbf{\arabic*.}, leftmargin=*, itemsep=0.28em]
  \item Si $f$ est continue sur $\intervalle{a}{b}$ et si $f(a)$ et $f(b)$
        sont de signes contraires, alors l'équation $f(x) = 0$ a une solution
        dans $\intervalle{a}{b}$.
  \item Une fonction admettant une limite finie en $a$ est continue en $a$.
  \item Si $u(x) \leq f(x) \leq v(x)$ avec $u \to 1$ et $v \to 1$, alors
        $f \to 1$.
  \item Toute fonction définie sur $\R$ est prolongeable par continuité.
  \item Si $f$ est continue et strictement décroissante sur
        $\intervalle{0}{5}$ avec $f(0) = 4$ et $f(5) = -2$, alors l'équation
        $f(x) = 3$ a exactement une solution.
\end{enumerate}
\end{jemeteste}

\begin{corrige}[je me teste]
\textbf{1.} Vrai — $0$ est compris entre $f(a)$ et $f(b)$.
\textbf{2.} Faux — encore faut-il que $f$ soit définie en $a$ et que la limite
vaille $f(a)$.
\textbf{3.} Vrai — c'est le théorème des gendarmes.
\textbf{4.} Faux — la question du prolongement ne se pose qu'aux points où la
fonction n'est pas définie.
\textbf{5.} Vrai — continuité, stricte monotonie, et $-2 < 3 < 4$.
\end{corrige}

% ===========================================================================
\begin{commeaubac}
\bmtsource{D'après le baccalauréat, série S}
On considère la fonction $f$ définie sur $\R \setminus \{-1\}$ par
\[
  f(x) = \frac{x^{2} + 2x + 3}{x + 1}.
\]

\textbf{Partie A}
\begin{enumerate}[label=\textbf{\arabic*.}, leftmargin=*, itemsep=0.25em]
  \item Déterminer les limites de $f$ en $-\infty$, en $+\infty$, puis à
        gauche et à droite de $-1$.
  \item Montrer que pour tout $x \neq -1$, $f(x) = x + 1 + \dfrac{2}{x+1}$.
  \item En déduire que la droite d'équation $y = x+1$ est asymptote à la
        courbe de $f$ en $+\infty$ et en $-\infty$.
\end{enumerate}

\textbf{Partie B}
\begin{enumerate}[label=\textbf{\arabic*.}, leftmargin=*, itemsep=0.25em, start=4]
  \item Calculer $f'(x)$ et dresser le tableau de variations de $f$.
  \item Montrer que l'équation $f(x) = 5$ admet exactement deux solutions
        réelles.
  \item Donner un encadrement d'amplitude $10^{-1}$ de la solution
        appartenant à $\left]-1 ; +\infty\right[$.
\end{enumerate}
\end{commeaubac}

\begin{corrige}[comme au baccalauréat]
\textbf{1.} En $\pm\infty$, $f(x) \sim x \to \pm\infty$. En $-1^{-}$ :
le numérateur tend vers $2 > 0$ et le dénominateur vers $0^{-}$, donc
$f \to -\infty$ ; en $-1^{+}$, $f \to +\infty$.

\textbf{2.} $x + 1 + \dfrac{2}{x+1} = \dfrac{(x+1)^{2} + 2}{x+1}
= \dfrac{x^{2}+2x+3}{x+1}$.

\textbf{3.} $f(x) - (x+1) = \dfrac{2}{x+1} \to 0$ en $\pm\infty$.

\textbf{4.} $f'(x) = 1 - \dfrac{2}{(x+1)^{2}} = \dfrac{(x+1)^{2}-2}{(x+1)^{2}}$,
qui s'annule pour $x = -1 \pm \sqrt{2}$. $f$ croît sur
$\left]-\infty ; -1-\sqrt{2}\right]$, décroît sur
$\left[-1-\sqrt{2} ; -1\right[$ puis sur $\left]-1 ; -1+\sqrt{2}\right]$, et
croît sur $\left[-1+\sqrt{2} ; +\infty\right[$. Les extremums valent
$f(-1\pm\sqrt{2}) = \pm 2\sqrt{2}$.

\textbf{5.} Sur $\left]-1 ; +\infty\right[$, le minimum est $2\sqrt{2} \approx 2{,}83 < 5$
et $f \to +\infty$ aux deux bords : deux branches, mais $f$ est strictement
monotone sur chacune. Sur $\left]-1 ; -1+\sqrt{2}\right]$, $f$ décroît de
$+\infty$ à $2\sqrt{2}$, donc $f(x)=5$ y a une solution unique. Sur
$\left[-1+\sqrt{2} ; +\infty\right[$, $f$ croît de $2\sqrt{2}$ à $+\infty$ :
une seconde solution unique. Sur $\left]-\infty ; -1\right[$, le maximum vaut
$-2\sqrt{2} < 5$ : aucune solution. Total : deux solutions.

\textbf{6.} $f(0{,}7) \approx 4{,}88$ et $f(0{,}8) \approx 5{,}02$, donc la
solution positive est comprise entre $0{,}7$ et $0{,}8$.
\end{corrige}

% ===========================================================================
\begin{concours}
\bmtsource{D'après le concours d'entrée de l'ENI, Fianarantsoa}
Soit $m$ un paramètre réel et $f_{m}$ la fonction définie par
\[
  f_{m}(x) = \frac{x^{2} + 3x + 4m}{x^{2} + (5m+1)x + 3}.
\]

\begin{enumerate}[label=\textbf{\arabic*.}, leftmargin=*, itemsep=0.25em]
  \item Déterminer $\lim_{x \to \pm\infty} f_{m}(x)$. Le résultat dépend-il
        de $m$ ?
  \item Montrer que toutes les courbes $\mathscr{C}_{m}$ passent par trois
        points fixes dont on déterminera les coordonnées.
  \item Pour $m = 0$, étudier la continuité de $f_{0}$ et préciser les valeurs
        où un prolongement par continuité est possible.
\end{enumerate}
\end{concours}

\begin{corrige}[vers le concours]
\textbf{1.} Les degrés du numérateur et du dénominateur sont égaux, de
coefficients dominants $1$ et $1$ : la limite vaut $1$ aux deux infinis, quelle
que soit la valeur de $m$.

\textbf{2.} Un point $(x_{0}, y_{0})$ est fixe lorsque l'égalité
$y_{0}\left(x_{0}^{2} + (5m+1)x_{0} + 3\right) = x_{0}^{2} + 3x_{0} + 4m$ vaut
pour tout $m$. En regroupant les termes en $m$ :
$m\left(5x_{0}y_{0} - 4\right) + \left(y_{0}x_{0}^{2} + y_{0}x_{0} + 3y_{0} - x_{0}^{2} - 3x_{0}\right) = 0$.
Il faut donc annuler séparément les deux parenthèses, ce qui conduit au
système à résoudre.

\textbf{3.} Pour $m = 0$, $f_{0}(x) = \dfrac{x^{2}+3x}{x^{2}+x+3}$. Le
discriminant du dénominateur vaut $1 - 12 = -11 < 0$ : il ne s'annule jamais.
La fonction est donc continue sur $\R$ tout entier et aucun prolongement n'a
lieu d'être.
\end{corrige}

% ===========================================================================
\begin{devoir}[2 heures]
Trois exercices indépendants, notés sur $20$. Le devoir porte sur ce chapitre
et sur les acquis antérieurs. Les énoncés sont repris de sessions réelles et
assemblés ici en épreuve ; leur provenance est indiquée à chaque fois.

\exods{1}{D'après le baccalauréat probatoire, Cameroun, série D, session 2024 — exercice III}{5 points}
\begin{enumerate}[label=\textbf{\arabic*.}, leftmargin=*, itemsep=0.2em]
  \item Démontrer que pour tout réel $x$, $-2+\cos x < 0$.
        \hfill\textup{(0,5 pt)}
  \item Démontrer que pour tout réel $x$,
        \[
          -3\cos x-2\sin^{2}x = \left(1+2\cos x\right)\left(-2+\cos x\right).
        \]
        \hfill\textup{(1 pt)}
  \item Résoudre dans $\intervallefo{0}{2\pi}$ l'équation
        $-3\cos x-2\sin^{2}x = 0$. \hfill\textup{(2 pts)}
  \item Résoudre dans $\intervallefo{0}{2\pi}$ l'inéquation
        $-3\cos x-2\sin^{2}x > 0$. \hfill\textup{(1,5 pt)}
\end{enumerate}

\exods{2}{D'après le baccalauréat probatoire, Cameroun, série C, session 2020 — exercice 3}{7 points}
Soit $f$ la fonction définie par $f(x) = ax+b+\dfrac{c}{x+1}$, où $a$, $b$ et
$c$ sont trois réels, et $(C)$ sa courbe. On sait que $(C)$ passe par les
points $A(0\,;6)$ et $B(-2\,;-4)$, et que la tangente à $(C)$ au point
$C(-3\,;-3)$ est parallèle à l'axe des abscisses.
\begin{enumerate}[label=\textbf{\arabic*.}, leftmargin=*, itemsep=0.2em]
  \item Déterminer les réels $a$, $b$ et $c$. \hfill\textup{(3 pts)}
  \item Donner l'ensemble de définition de $f$, puis calculer ses limites aux
        bornes de cet ensemble. \hfill\textup{(2 pts)}
  \item En déduire les asymptotes à la courbe $(C)$. \hfill\textup{(2 pts)}
\end{enumerate}

\exods{3}{D'après le baccalauréat probatoire, Cameroun, série D, session 2024 — exercice II}{8 points}
Soit $f$ la fonction définie sur $\R$ par $f(x) = -x^{3}+3x^{2}$, et $g$ la
fonction définie sur $\R$ par $g(x) = ax^{3}+bx^{2}+cx+4$, de courbe
$(\mathcal C_{g})$.
\begin{enumerate}[label=\textbf{\arabic*.}, leftmargin=*, itemsep=0.2em]
  \item Calculer les limites de $f$ en $-\infty$ et en $+\infty$.
        \hfill\textup{(1 pt)}
  \item Résoudre dans $\R$ l'équation $f(x) = 0$, puis donner le signe de
        $f(x)$ suivant les valeurs de $x$. \hfill\textup{(1,5 pt)}
  \item Sachant que $(\mathcal C_{g})$ coupe l'axe des abscisses au point
        d'abscisse $1$ et qu'elle admet au point $S(3\,;4)$ une tangente
        parallèle à l'axe des abscisses, déterminer $a$, $b$ et $c$.
        \hfill\textup{(3 pts)}
  \item Vérifier que $g(x) = f(x-1)$ pour tout réel $x$, et en déduire les
        limites de $g$ en $-\infty$ et en $+\infty$. \hfill\textup{(1,5 pt)}
  \item Résoudre l'équation $g(x) = 0$. \hfill\textup{(1 pt)}
\end{enumerate}
\end{devoir}

\begin{corrige}[devoir surveillé]
\textbf{Exercice 1. 1.} Pour tout réel $x$, $\cos x \leq 1 < 2$, donc
$-2+\cos x \leq -1 < 0$ : la quantité est strictement négative, et ne s'annule
jamais.

\textbf{2.} On développe le produit et l'on remplace $\sin^{2}x$ par
$1-\cos^{2}x$ :
\[
  (1+2\cos x)(-2+\cos x) = -2+\cos x-4\cos x+2\cos^{2}x
  = 2\cos^{2}x-3\cos x-2,
\]
tandis que
$-3\cos x-2\sin^{2}x = -3\cos x-2\left(1-\cos^{2}x\right)
= 2\cos^{2}x-3\cos x-2$. Les deux écritures coïncident.

\textbf{3.} Le produit est nul si l'un des facteurs l'est. Le second ne
s'annule jamais, d'après la question \textbf{1}. Reste
$1+2\cos x = 0$, soit $\cos x = -\frac12$, dont les solutions dans
$\intervallefo{0}{2\pi}$ sont $\dfrac{2\pi}{3}$ et $\dfrac{4\pi}{3}$.

\textbf{4.} Le second facteur étant strictement négatif, le produit est
strictement positif exactement lorsque le premier est strictement négatif :
$1+2\cos x < 0$, c'est-à-dire $\cos x < -\frac12$. Sur
$\intervallefo{0}{2\pi}$, cela donne
$x \in \left]\dfrac{2\pi}{3}\,;\dfrac{4\pi}{3}\right[$.

\medskip
\textbf{Exercice 2. 1.} Trois conditions, trois équations. Le passage par $A$
donne $b+c = 6$ ; le passage par $B$ donne $-2a+b-c = -4$ ; le passage par
$C$ donne $-3a+b-\frac{c}{2} = -3$. Enfin
$f'(x) = a-\dfrac{c}{(x+1)^{2}}$, et la tangente en $C$ étant horizontale,
$f'(-3) = a-\dfrac{c}{4} = 0$, soit $c = 4a$. En reportant dans la première
équation, $b = 6-4a$ ; dans la troisième,
$-3a+6-4a-2a = -3$, d'où $-9a = -9$ et $a = 1$. Ainsi
\[
  a = 1, \qquad b = 2, \qquad c = 4,
  \qquad\text{c'est-à-dire}\qquad
  f(x) = x+2+\frac{4}{x+1}.
\]
La deuxième équation est bien vérifiée : $-2+2-4 = -4$.

\textbf{2.} Le dénominateur s'annule en $-1$ : $\Df = \R\setminus\{-1\}$. Aux
infinis, $\frac{4}{x+1}\to0$ et $f(x)$ se comporte comme $x+2$ : $f(x)\to
-\infty$ en $-\infty$ et $f(x)\to+\infty$ en $+\infty$. En $-1^{-}$,
$\frac{4}{x+1}\to-\infty$ donc $f(x)\to-\infty$ ; en $-1^{+}$,
$f(x)\to+\infty$.

\textbf{3.} La droite $x = -1$ est asymptote verticale. Comme
$f(x)-(x+2) = \dfrac{4}{x+1}$ tend vers $0$ aux deux infinis, la droite
d'équation $y = x+2$ est asymptote oblique, la courbe étant au-dessus pour
$x>-1$ et au-dessous pour $x<-1$.

\medskip
\textbf{Exercice 3. 1.} $f(x) = x^{3}\left(-1+\frac{3}{x}\right)$ : en
$+\infty$ la parenthèse tend vers $-1$ et $x^{3}$ vers $+\infty$, donc
$f(x)\to-\infty$ ; en $-\infty$, $f(x)\to+\infty$.

\textbf{2.} $f(x) = x^{2}(3-x)$, nul pour $x = 0$ — racine double — et pour
$x = 3$. Comme $x^{2}\geq0$, le signe de $f$ est celui de $3-x$ : $f$ est
positive sur $\left]-\infty\,;3\right]$, négative sur
$\intervallefo{3}{+\infty}$, et s'annule en $0$ sans changer de signe.

\textbf{3.} La courbe coupe l'axe des abscisses en $1$ : $g(1) = 0$, soit
$a+b+c+4 = 0$. Le point $S(3\,;4)$ appartient à la courbe :
$27a+9b+3c+4 = 4$, soit $9a+3b+c = 0$. La tangente y étant horizontale,
$g'(3) = 27a+6b+c = 0$. La différence des deux dernières égalités donne
$-18a-3b = 0$, d'où $b = -6a$ ; en reportant, $c = 9a$ ; la première équation
devient $a-6a+9a+4 = 0$, soit $4a = -4$. Finalement
\[
  a = -1, \qquad b = 6, \qquad c = -9,
  \qquad\text{c'est-à-dire}\qquad
  g(x) = -x^{3}+6x^{2}-9x+4.
\]

\textbf{4.} On développe :
$f(x-1) = -(x-1)^{3}+3(x-1)^{2}
= -\left(x^{3}-3x^{2}+3x-1\right)+3\left(x^{2}-2x+1\right)
= -x^{3}+6x^{2}-9x+4 = g(x)$. La courbe de $g$ est donc l'image de celle de
$f$ par la translation qui déplace tout point d'une unité vers la droite ;
les limites sont inchangées : $g(x)\to+\infty$ en $-\infty$ et
$g(x)\to-\infty$ en $+\infty$.

\textbf{5.} $g(x) = f(x-1) = (x-1)^{2}\left(3-(x-1)\right)
= (x-1)^{2}(4-x)$. L'équation $g(x) = 0$ a donc pour solutions $x = 1$, racine
double, et $x = 4$.
\end{corrige}
